The uplink of the massive MIMO system was investigated using two types of parameters. Hard constraints refer to the physical characteristics of the base station: $M = 128$ or 256 antennas, $C = 8$ or 16 clusters, and the size $M_c = M/C$ of the inversion module. Soft constraints depend on the UEs being served and the level of interference: $K = 8$ or 12 UEs, QPSK or 16-QAM modulation, and an IoTh ratio of 10 dB as a baseline case, or 20 dB representing an extreme, high-interference scenario. These configurations were adapted from \cite{zhao2025jsac}. Table~\ref{tab:paineis} presents these configurations.


Each configuration is identified by a label consisting of a number and a letter. The number denotes the hard constraints: 1 for $M = 128$ and $C = 8$; 2 for $M = 256$ and $C = 8$; and 3 for $M = 256$ and $C = 16$. The letter indicates the soft constraints and remains consistent across all three groups. \textcolor{red}{As configurações com o mesmo número formam um grupo: o grupo 1 tem as configurações 1A--1H, o grupo 2, as 2A--2H, e o grupo 3, as 3A--3H.}

\begin{table}[htbp]
\caption{\textcolor{red}{Configurações simuladas.}}
\label{tab:paineis}
\vspace{4pt} % a classe zera o espaço entre a caption e a tabela
\centering
\footnotesize
\renewcommand{\arraystretch}{1.2}
% sem espaço em volta das linhas do booktabs, para as bordas verticais não terem falhas
\setlength{\aboverulesep}{0pt}
\setlength{\belowrulesep}{0pt}
\begin{tabular}{|c|c|c|c|c|c|c|}
\toprule
\multirow{2}{*}{\textbf{Configuration}} & \multicolumn{3}{c|}{\textcolor{red}{\textbf{Hard constraints}}} & \multicolumn{3}{c|}{\textcolor{red}{\textbf{Soft constraints}}} \\
\cline{2-7}
 & {\boldmath $M$} & {\boldmath $C$} & {\boldmath $M_c$} & \textbf{IoTh (dB)} & \textbf{Modulation} & {\boldmath $K$} \\
-- & \texttt{M} & \texttt{C} & -- & \texttt{iot\_db} & \texttt{modulation} & \texttt{K} \\
\midrule
1A & \multirow{8}{*}{128} & \multirow{8}{*}{8} & \multirow{8}{*}{16} & \multirow{4}{*}{10} & \multirow{2}{*}{16-QAM} & 8 \\
1B & & & & & & 12 \\
\cline{6-7}
1C & & & & & \multirow{2}{*}{QPSK} & 8 \\
1D & & & & & & 12 \\
\cline{5-7}
1E & & & & \multirow{4}{*}{20} & \multirow{2}{*}{16-QAM} & 8 \\
1F & & & & & & 12 \\
\cline{6-7}
1G & & & & & \multirow{2}{*}{QPSK} & 8 \\
1H & & & & & & 12 \\
\midrule
2A & \multirow{8}{*}{256} & \multirow{8}{*}{8} & \multirow{8}{*}{32} & \multirow{4}{*}{10} & \multirow{2}{*}{16-QAM} & 8 \\
2B & & & & & & 12 \\
\cline{6-7}
2C & & & & & \multirow{2}{*}{QPSK} & 8 \\
2D & & & & & & 12 \\
\cline{5-7}
2E & & & & \multirow{4}{*}{20} & \multirow{2}{*}{16-QAM} & 8 \\
2F & & & & & & 12 \\
\cline{6-7}
2G & & & & & \multirow{2}{*}{QPSK} & 8 \\
2H & & & & & & 12 \\
\midrule
3A & \multirow{8}{*}{256} & \multirow{8}{*}{16} & \multirow{8}{*}{16} & \multirow{4}{*}{10} & \multirow{2}{*}{16-QAM} & 8 \\
3B & & & & & & 12 \\
\cline{6-7}
3C & & & & & \multirow{2}{*}{QPSK} & 8 \\
3D & & & & & & 12 \\
\cline{5-7}
3E & & & & \multirow{4}{*}{20} & \multirow{2}{*}{16-QAM} & 8 \\
3F & & & & & & 12 \\
\cline{6-7}
3G & & & & & \multirow{2}{*}{QPSK} & 8 \\
3H & & & & & & 12 \\
\bottomrule
\end{tabular}
\vspace{4pt}
\legend{\textcolor{red}{Fonte: configurações 1A, 1B, 1C, 1E, 2A e 3A adaptadas de \citesp{zhao2025jsac}; as outras, a autora. Em todas as configurações, $r = 8$. Nomes no código: campos de \texttt{Panel} (\texttt{dbp\_sim/params.py}).}}
% EN (versão anterior, 2026-11-05): Source: configurations 1A, 1B, 1C, 1E, 2A and 3A adapted from \citesp{zhao2025jsac}; the others, the author. In all configurations, $r = 8$. Names in the code: fields of \texttt{Panel} (\texttt{dbp\_sim/params.py}).
\end{table}

% AUTORA (9-results.tex:245) -- parágrafo retirado do texto (tom); o conteúdo ficou em comentário.
% Configurações 1, 9, 13, 3, 2, 5 = painéis (a), (b), (c), (e), (f), (g) da Fig. 8 de Zhao et al.
% O painel (d) não aparece no artigo. Na configuração 3 (painel e), r = 8 foi
% mantido; L577 liga r a K, então r = 12 também seria defensável.


